% GENERATED FILE. Edit canonical lessons, then run npm run generate:books.
% canonical-source-hash: fnv1a64:ad5992531b4605ae
% canonical-lessons: GU-C362-sansodhan, GU-C362-prayog, GU-C362-paddhati, GU-C362-namuno, GU-C362-taravvu, GU-R362-review-355-362

\chapter{Research}
\label{ch:gu-research}
\hlchaptermodality{362}

\begin{chapteropening}
\textbf{By the end of this chapter:} \emph{I can say what a researcher is working on.}

Complete the last lesson of chapter 362: I can say what a researcher is working on.
\end{chapteropening}

\section[\texorpdfstring{\gu{સંશોધન} (sanśodhan)}{સંશોધન (sanśodhan)}]{\gu{સંશોધન} (sanśodhan) — research}

\label{lesson:GU-C362-sansodhan}



\begin{quote}
Before the new one: say the Gujarati for to run down, then the Gujarati for to entreat.
\end{quote}

\subsection*{You'll want to know: \gu{સંશોધન}}
\textbf{\gu{સંશોધન}} — \emph{sanśodhan} — \textquotedblleft{}research\textquotedblright{}.

\begin{quote}
\gu{મીરા} \gu{નવી} \gu{દવા} \gu{પર} \gu{સંશોધન} \gu{કરે} \gu{છે}.
\end{quote}

\begin{quote}
\emph{Mīrā navī davā par sanśodhan kare che.}
\end{quote}

\begin{quote}
Mira is doing research on a new medicine.
\end{quote}

A neuter noun.

\begin{cousinweb}
\textbf{\gu{સંશોધન}} (\emph{sanśodhan}) is a tatsama, borrowed from Sanskrit \emph{saṃśodhana}, \textquotedblleft{}thoroughly cleansing\textquotedblright{} (Monier-Williams); \textquotedblleft{}research\textquotedblright{}, sifting a question clean, is its modern sense.
\end{cousinweb}

\begin{grammarlens}[title={What you've built}]
One more word for giving reasons and opinions.
\end{grammarlens}

\subsection*{Your turn}
\begin{itemize}
\raggedright
  \item \emph{Say it:} \emph{sanśodhan}
  \item \emph{Say it:} \emph{sanśodhan}, once more
  \item \emph{Say it:} \emph{Mīrā navī davā par sanśodhan kare che.} — Mira is doing research on a new medicine.
  \item \emph{Recall:} say the Gujarati for to run down, then the Gujarati for to entreat, then say \emph{sanśodhan} again
\end{itemize}

\begin{tcolorbox}[breakable,colback=teal!4,colframe=teal!35!black,title={Before you move on}]
What does \gu{સંશોધન} mean? (\textquotedblleft{}Research\textquotedblright{}.) Say it once more.
\end{tcolorbox}
\section[\texorpdfstring{\gu{પ્રયોગ} (prayog)}{પ્રયોગ (prayog)}]{\gu{પ્રયોગ} (prayog) — an experiment}

\label{lesson:GU-C362-prayog}



\begin{quote}
Before the new one: say the Gujarati for to entreat, then the Gujarati for research.
\end{quote}

\subsection*{You'll want to know: \gu{પ્રયોગ}}
\textbf{\gu{પ્રયોગ}} — \emph{prayog} — \textquotedblleft{}an experiment\textquotedblright{}.

\begin{quote}
\gu{વિદ્યાર્થીઓએ} \gu{વર્ગમાં} \gu{પાણી} \gu{પર} \gu{એક} \gu{પ્રયોગ} \gu{કર્યો}.
\end{quote}

\begin{quote}
\emph{vidyārthīoe vargamā̃ pāṇī par ek prayog karyo.}
\end{quote}

\begin{quote}
The students did an experiment with water in class.
\end{quote}

A masculine noun.

\begin{cousinweb}
\textbf{\gu{પ્રયોગ}} (\emph{prayog}) is a tatsama, borrowed from Sanskrit \emph{prayoga}, \textquotedblleft{}use, application\textquotedblright{} (Monier-Williams): an idea put to use.
\end{cousinweb}

\begin{grammarlens}[title={What you've built}]
One more word for giving reasons and opinions.
\end{grammarlens}

\subsection*{Your turn}
\begin{itemize}
\raggedright
  \item \emph{Say it:} \emph{prayog}
  \item \emph{Say it:} \emph{prayog}, once more
  \item \emph{Say it:} \emph{vidyārthīoe vargamā̃ pāṇī par ek prayog karyo.} — The students did an experiment with water in class.
  \item \emph{Recall:} say the Gujarati for to entreat, then the Gujarati for research, then say \emph{prayog} again
\end{itemize}

\begin{tcolorbox}[breakable,colback=teal!4,colframe=teal!35!black,title={Before you move on}]
What does \gu{પ્રયોગ} mean? (\textquotedblleft{}An experiment\textquotedblright{}.) Say it once more.
\end{tcolorbox}
\section[\texorpdfstring{\gu{પદ્ધતિ} (paddhati)}{પદ્ધતિ (paddhati)}]{\gu{પદ્ધતિ} (paddhati) — a method}

\label{lesson:GU-C362-paddhati}



\begin{quote}
Before the new one: say the Gujarati for research, then the Gujarati for an experiment.
\end{quote}

\subsection*{You'll want to know: \gu{પદ્ધતિ}}
\textbf{\gu{પદ્ધતિ}} — \emph{paddhati} — \textquotedblleft{}a method\textquotedblright{}.

\begin{quote}
\gu{આ} \gu{પદ્ધતિથી} \gu{કામ} \gu{ઝડપી} \gu{બને} \gu{છે}.
\end{quote}

\begin{quote}
\emph{ā paddhatithī kām jhaḍpī bane che.}
\end{quote}

\begin{quote}
This method makes the work go faster.
\end{quote}

A feminine noun.

\begin{cousinweb}
\textbf{\gu{પદ્ધતિ}} (\emph{paddhati}) is a tatsama, borrowed from Sanskrit \emph{paddhati}, \textquotedblleft{}a path, a way, a course\textquotedblright{} (Monier-Williams).
\end{cousinweb}

\begin{grammarlens}[title={What you've built}]
One more word for giving reasons and opinions.
\end{grammarlens}

\subsection*{Your turn}
\begin{itemize}
\raggedright
  \item \emph{Say it:} \emph{paddhati}
  \item \emph{Say it:} \emph{paddhati}, once more
  \item \emph{Say it:} \emph{ā paddhatithī kām jhaḍpī bane che.} — This method makes the work go faster.
  \item \emph{Recall:} say the Gujarati for research, then the Gujarati for an experiment, then say \emph{paddhati} again
\end{itemize}

\begin{tcolorbox}[breakable,colback=teal!4,colframe=teal!35!black,title={Before you move on}]
What does \gu{પદ્ધતિ} mean? (\textquotedblleft{}A method\textquotedblright{}.) Say it once more.
\end{tcolorbox}
\section[\texorpdfstring{\gu{નમૂનો} (namūno)}{નમૂનો (namūno)}]{\gu{નમૂનો} (namūno) — a sample}

\label{lesson:GU-C362-namuno}



\begin{quote}
Before the new one: say the Gujarati for an experiment, then the Gujarati for a method.
\end{quote}

\subsection*{You'll want to know: \gu{નમૂનો}}
\textbf{\gu{નમૂનો}} — \emph{namūno} — \textquotedblleft{}a sample\textquotedblright{}.

\begin{quote}
\gu{દાક્તરે} \gu{લોહીનો} \gu{નમૂનો} \gu{લીધો}.
\end{quote}

\begin{quote}
\emph{dāktare lohīno namūno līdho.}
\end{quote}

\begin{quote}
The doctor took a blood sample.
\end{quote}

A masculine noun.

\begin{cousinweb}
\textbf{\gu{નમૂનો}} (\emph{namūno}) is usually traced to Persian \emph{namūna}, \textquotedblleft{}a pattern, a sample\textquotedblright{}.
\end{cousinweb}

\begin{grammarlens}[title={What you've built}]
One more word for giving reasons and opinions.
\end{grammarlens}

\subsection*{Your turn}
\begin{itemize}
\raggedright
  \item \emph{Say it:} \emph{namūno}
  \item \emph{Say it:} \emph{namūno}, once more
  \item \emph{Say it:} \emph{dāktare lohīno namūno līdho.} — The doctor took a blood sample.
  \item \emph{Recall:} say the Gujarati for an experiment, then the Gujarati for a method, then say \emph{namūno} again
\end{itemize}

\begin{tcolorbox}[breakable,colback=teal!4,colframe=teal!35!black,title={Before you move on}]
What does \gu{નમૂનો} mean? (\textquotedblleft{}A sample\textquotedblright{}.) Say it once more.
\end{tcolorbox}
\section[\texorpdfstring{\gu{તારવવું} (tāravvũ)}{તારવવું (tāravvũ)}]{\gu{તારવવું} (tāravvũ) — to deduce}

\label{lesson:GU-C362-taravvu}



\begin{quote}
Before the new one: say the Gujarati for a method, then the Gujarati for a sample.
\end{quote}

\subsection*{You'll want to know: \gu{તારવવું}}
\textbf{\gu{તારવવું}} — \emph{tāravvũ} — \textquotedblleft{}to deduce\textquotedblright{}.

\begin{quote}
\gu{આ} \gu{પ્રયોગમાંથી} \gu{મીરાએ} \gu{એક} \gu{વાત} \gu{તારવી}.
\end{quote}

\begin{quote}
\emph{ā prayogmā̃thī Mīrāe ek vāt tārvī.}
\end{quote}

\begin{quote}
From this experiment Mira drew one conclusion.
\end{quote}

Past, with \gu{તેણે} (he or she …): \gu{તેણે} \gu{સાચો} \gu{અર્થ} \gu{તારવ્યો}, \emph{teṇe sāco arth tāravyo}.

\begin{cousinweb}
\textbf{\gu{તારવવું}} (\emph{tāravvũ}) is inherited from Sanskrit \emph{tāráyati}, \textquotedblleft{}causes to arrive at, leads across\textquotedblright{}; CDIAL lists the Gujarati verb, \textquotedblleft{}to raise, to pluck\textquotedblright{} (CDIAL 5796). \textquotedblleft{}To deduce\textquotedblright{}, to pluck a conclusion out of the facts, is not among CDIAL's senses.
\end{cousinweb}

\begin{grammarlens}[title={What you've built}]
One more word for giving reasons and opinions.
\end{grammarlens}

\subsection*{Your turn}
\begin{itemize}
\raggedright
  \item \emph{Say it:} \emph{tāravvũ}
  \item \emph{Say it:} \emph{tāravvũ}, once more
  \item \emph{Say it:} \emph{ā prayogmā̃thī Mīrāe ek vāt tārvī.} — From this experiment Mira drew one conclusion.
  \item \emph{Recall:} say the Gujarati for a method, then the Gujarati for a sample, then say \emph{tāravvũ} again
\end{itemize}

\begin{tcolorbox}[breakable,colback=teal!4,colframe=teal!35!black,title={Before you move on}]
What does \gu{તારવવું} mean? (\textquotedblleft{}To deduce\textquotedblright{}.) Say it once more.
\end{tcolorbox}
\section[(dialogue)]{Looking back: from Taxes and Public Money to Research}

\label{lesson:GU-R362-review-355-362}



\begin{quote}
Say the Gujarati for a tax and a contribution.
\end{quote}

\subsection*{Your turn}
Say these aloud:

\begin{itemize}
\raggedright
  \item a tax, a contribution, economic, a grant, to collect money
  \item a constitution, independent, society, valid, to preserve
  \item a moral value, upbringing, a policy, a social practice, to adopt
  \item service to others, giving something up, doing good to others, cooperation, to keep up
  \item one's good name, dignity, shame, conceit, to bow
  \item integrity, temptation, loyal, a betrayal, to entrap
  \item speaking ill of someone, criticism, a telling-off, to run down, to entreat
  \item research, an experiment, a method, a sample, to deduce
\end{itemize}

\begin{tcolorbox}[breakable,colback=teal!4,colframe=teal!35!black,title={Before you move on}]
A tax? (\textbf{\gu{વેરો}}, \emph{vero}.) A contribution? (\textbf{\gu{ફાળો}}, \emph{phāḷo}.) Economic? (\textbf{\gu{આર્થિક}}, \emph{ārthik}.) A grant? (\textbf{\gu{ગ્રાન્ટ}}, \emph{grānṭ}.) To collect money? (\textbf{\gu{ઉઘરાવવું}}, \emph{ughrāvvũ}.) A constitution? (\textbf{\gu{બંધારણ}}, \emph{bandhāraṇ}.) Independent? (\textbf{\gu{સ્વતંત્ર}}, \emph{svatantra}.) Society? (\textbf{\gu{સમાજ}}, \emph{samāj}.) Valid? (\textbf{\gu{માન્ય}}, \emph{mānya}.) To preserve? (\textbf{\gu{જાળવવું}}, \emph{jāḷavvũ}.) A moral value? (\textbf{\gu{મૂલ્ય}}, \emph{mūlya}.) Upbringing? (\textbf{\gu{સંસ્કાર}}, \emph{sanskār}.) A policy? (\textbf{\gu{નીતિ}}, \emph{nīti}.) A social practice? (\textbf{\gu{પ્રથા}}, \emph{prathā}.) To adopt? (\textbf{\gu{અપનાવવું}}, \emph{apnāvvũ}.) Service to others? (\textbf{\gu{સેવા}}, \emph{sevā}.) Giving something up? (\textbf{\gu{ત્યાગ}}, \emph{tyāg}.) Doing good to others? (\textbf{\gu{પરોપકાર}}, \emph{paropkār}.) Cooperation? (\textbf{\gu{સહકાર}}, \emph{sahkār}.) To keep up? (\textbf{\gu{નિભાવવું}}, \emph{nibhāvvũ}.) One's good name? (\textbf{\gu{ઇજ્જત}}, \emph{ijjat}.) Dignity? (\textbf{\gu{ગૌરવ}}, \emph{gaurav}.) Shame? (\textbf{\gu{શરમ}}, \emph{śaram}.) Conceit? (\textbf{\gu{અભિમાન}}, \emph{abhimān}.) To bow? (\textbf{\gu{નમવું}}, \emph{namvũ}.) Integrity? (\textbf{\gu{ઈમાન}}, \emph{īmān}.) Temptation? (\textbf{\gu{લાલચ}}, \emph{lālac}.) Loyal? (\textbf{\gu{વફાદાર}}, \emph{vaphādār}.) A betrayal? (\textbf{\gu{દગો}}, \emph{dago}.) To entrap? (\textbf{\gu{ફસાવવું}}, \emph{phasāvvũ}.) Speaking ill of someone? (\textbf{\gu{નિંદા}}, \emph{nindā}.) Criticism? (\textbf{\gu{ટીકા}}, \emph{ṭīkā}.) A telling-off? (\textbf{\gu{ઠપકો}}, \emph{ṭhapko}.) To run down? (\textbf{\gu{વગોવવું}}, \emph{vagovvũ}.) To entreat? (\textbf{\gu{વીનવવું}}, \emph{vīnavvũ}.) Research? (\textbf{\gu{સંશોધન}}, \emph{sanśodhan}.) An experiment? (\textbf{\gu{પ્રયોગ}}, \emph{prayog}.) A method? (\textbf{\gu{પદ્ધતિ}}, \emph{paddhati}.) A sample? (\textbf{\gu{નમૂનો}}, \emph{namūno}.) To deduce? (\textbf{\gu{તારવવું}}, \emph{tāravvũ}.)
\end{tcolorbox}
